\section{Complete proof of the global upper bound}
\label{app:upper}

This appendix defines the estimator and proves every stochastic and
deterministic step used in Proposition~\ref{prop:upper}.  No result from an
internal proof package is an assumption below.  Throughout, \(P_j\) denotes
the empirical average on half \(\cD_j\), \(j=1,2\).  Every unnamed proof
constant depends only on the fixed tuple \(\rho\).  Throughout, put
\(v:=\E u^2\in[v_-,v_+]\); no constant or positive pointwise
conditional treatment variance is assumed.  Put
\[
 \norm{f}_{1,n}:=(P_1f^2)^{1/2},\qquad
 \norm{f}_{2,n}:=(P_2f^2)^{1/2}.
\]
These half-indexed symbols are used for fitted functions of \(X\).  A vector
norm whose sample half is already fixed by its definition retains the shorter
notation \(\norm\cdot_n\).
Put
\[
 C_f=1+\frac{B_0}{\sqrt{v_-}},\qquad
 C_J=1+3C_f+\frac{B_0}{c_{\rm rad}\sqrt{v_-}}.
\]
Let \(K_{1,\star}\) be the maximum of the primitive radius constants in
Lemma~\ref{lem:localized-process} for the seven calls listed explicitly in
the event ledger below.  They use only the two public difference cores and
the joint class, on a specified sample half, with a specified multiplier or
square process, before any fitted index is substituted.  Fix
\[
 C_{\rm rad}:=1+K_{1,\star}(1+C_J),\qquad
 K_\lambda:=1,\qquad K_d:=8.
\]
Here the model parameter \(c_{\rm rad}>0\) is distinct from the
proof/estimator inflation constant \(C_{\rm rad}\).  Also put
\[
 B_\pi=B_0+B_{\mathcal G},\qquad
 B_{\rm all}=B_0+4B_{\mathcal G},\qquad
 \vartheta_\rho=\frac{B_\pi}{8},\qquad
 \delta_\rho=\frac{\vartheta_\rho}{C_{\rm rad}}.
\]
These quantities are fixed before \(n,(a,s,b,t),\Gamma,P\), and the data;
they are never enlarged or redefined by a later proof constant.  Finally set
\[
 h_n=\sqrt{\frac{\log(en)}n},\quad
 \bar s=C_{\rm rad}(s\vee h_n),\quad
 \bar t=C_{\rm rad}(t\vee h_n),\quad
 \rho_n=n^{-6},\quad \tau_n=n^{-12}.
\]

\subsection{Elementary concentration and countable optimization}

\begin{lemma}[Bounded Bernstein inequality]
\label{lem:bernstein}
Let \(\mathcal F\) be a sigma-field.  Suppose \(W_1,\ldots,W_n\) are
conditionally independent given \(\mathcal F\) and conditionally centered
given \(\mathcal F\), with \(\abs{W_i}\le B\), and
\(n^{-1}\sum_i\E(W_i^2\mid\mathcal F)\le\sigma^2\), then, conditionally on
\(\mathcal F\), with probability at least \(1-2e^{-z}\),
\[
 \abs{P_nW}\le \sqrt{\frac{2\sigma^2z}{n}}+\frac{2Bz}{3n}.
 \tag{B.1}
\]
\end{lemma}
\begin{proof}
For \(0<\lambda<3/B\), the power series and
\(k!\ge2\,3^{k-2}\) for \(k\ge2\) give
\[
 \E(e^{\lambda W_i}\mid\mathcal F)
 \le1+\frac{\lambda^2\E(W_i^2\mid\mathcal F)}
 {2(1-\lambda B/3)}.
\]
Use \(1+x\le e^x\).  Conditional independence given \(\mathcal F\) permits
multiplication of the conditional moment generating functions.  Apply
Markov's inequality and optimize in \(\lambda\).  This gives the
upper tail in (B.1); apply the same argument to \(-W_i\) for the lower tail.
\end{proof}

We also use the following weighted form, which does not require a bound on
the largest weight.  If \(c_1,\ldots,c_n\) are \(\mathcal F\)-measurable and,
conditionally on \(\mathcal F\), the \(\xi_i\) are conditionally independent
given \(\mathcal F\), conditionally centered given \(\mathcal F\), and satisfy
\(\abs{\xi_i}\le B\), then
\[
 \abs{P_n(\xi c)}
 \le B\norm c_n\sqrt{\frac{2z}{n}}
 \tag{B.1a}
\]
with conditional failure at most \(2e^{-z}\).  Indeed, Hoeffding's
one-variable exponential bound and conditional independence give
\(\E[\exp\{\lambda\sum_i c_i\xi_i\}\mid\mathcal F]
\le\exp\{\lambda^2B^2\sum_i c_i^2/2\}\); optimize the two Chernoff tails.

\begin{lemma}[Localized supremum concentration]
\label{lem:supremum-bernstein}
Let \(\mathcal W\) be a countable class of centered functions of an iid
observation, with \(\abs{W(w)}\le B\) and
\(\E W(w)^2\le\sigma^2\), uniformly in \(w\).  Put
\[
 \mathcal R_n(\mathcal W)
 =\E_{\mathrm{data},\epsilon}\sup_{w\in\mathcal W}
   \abs{\frac1n\sum_{i=1}^n\epsilon_iW_i(w)}.
\]
If \(Z=\sup_{w\in\mathcal W}\abs{P_nW(w)}\), then
\[
 \Pp\!\left\{Z>4\mathcal R_n(\mathcal W)
            +\sqrt{\frac{2\sigma^2z}{n}}
            +\frac{8Bz}{3n}\right\}\le2e^{-z}.
 \tag{B.2}
\]
\end{lemma}
\begin{proof}
This is an exact transfer of Theorem 2.1 of
\citet[pp. 1503--1504]{bartlett2005local}.  Enlarge \(\mathcal W\) by its
negatives and map the source theorem's \([a,b]\), variance parameter \(r\),
confidence \(x\), and tuning \(\alpha\) to
\([-B,B]\), \(\sigma^2\), \(z\), and \(1\), respectively.  Its one-sided
supremum is then exactly \(Z\); its Rademacher term is
\(\mathcal R_n(\mathcal W)\); and its range term is
\(2B(1/3+1)z/n=8Bz/(3n)\).  The source probability is \(1-e^{-z}\), so
the displayed \(2e^{-z}\) is conservative.  The source assumes measurable
suprema; countability supplies this here.  Thus every source assumption
has been mapped, and no project-internal concentration certificate is used.
\end{proof}

\begin{lemma}[Localized multiplier and square processes]
\label{lem:localized-process}
Let \(\mathcal H\) be a countable class with envelope \(B\), containing zero,
and suppose
\[
 \mathfrak R_{n,P_X}(r;\mathcal H)\le c_{\rm rad}r_0r,
 \qquad r\ge r_0,\quad r>0.
 \tag{B.2a}
\]
Let \(\xi\) be bounded and satisfy \(\E(\xi\mid X)=0\).  For
\(\delta_z=K_1(r_0\vee h_n)+\sqrt{z/n}\) and \(1\le z\le n\), an event
of probability at least \(1-K_2e^{-z}\) satisfies, simultaneously for every
\(f\in\mathcal H\),
\begin{align}
 \abs{(P_n-P)(\xi f)}
 &\le K_3\{\delta_z\norm f_2+\delta_z^2\},\tag{B.3}\\
 \abs{(P_n-P)f^2}
 &\le K_3\{\delta_z\norm f_2+\delta_z^2\}.\tag{B.4}
\end{align}
The same conclusion holds when \(\xi\) is a fixed bounded multiplier, with
the centered process in (B.3).
\end{lemma}
\begin{proof}
Write \(L=\norm\xi_\infty\).  Define the inner ball and dyadic shells by
\[
 \mathcal H_0=\{f:\norm f_2\le\delta_z\},\qquad
 \mathcal H_j=\{f:2^{j-1}\delta_z<\norm f_2
       \le\min(2^j\delta_z,B)\},\quad j\ge1,
\]
and set
\[
 J=\left\lceil\log_2(B/\delta_z)\right\rceil_+.
\]
The envelope gives \(\mathcal H=\bigcup_{j=0}^J\mathcal H_j\); if
\(\delta_z\ge B\), then \(J=0\).  Since \(\delta_z\ge h_n\),
\begin{equation}
 J\le \left\lceil\log_2(B/h_n)\right\rceil_+
 \le C_B\log(en),\qquad
 \sqrt{j/n}\le C_Bh_n,\qquad j/n\le C_Bh_n^2.
 \tag{B.2b}
\end{equation}

For \(j\ge1\), put \(R_j=2^j\delta_z\), and put
\(R_0=\delta_z\).  On \(\mathcal H_j\), conditional Rademacher
contraction bounds the multiplier mean by
\(L\mathfrak R_n(R_j;\mathcal H)\).  If the multiplier is fixed rather
than conditionally centered, centering adds at most \(LR_j/\sqrt n\).
Thus its centered Rademacher mean is at most
\(C(r_0+h_n)R_j\le C\delta_zR_j\), its variance is at most
\(L^2R_j^2\), and its centered range is at most \(2LB\).
These multiplier bounds use only boundedness, not a constant conditional
variance.

For the square process,
\[
 \operatorname{Var}(f^2)\le \E f^4\le B^2R_j^2.
\]
The map \(u\mapsto u^2\) is \(2B\)-Lipschitz on \([-B,B]\), so
symmetrization and contraction bound the uncentered Rademacher term by
\(CB\mathfrak R_n(R_j;\mathcal H)\).  By localization and the envelope,
the centering contribution satisfies
\[
 \sup_{f\in\mathcal H_j}Pf^2\,\E_\epsilon
 \abs{\frac1n\sum_{i=1}^n\epsilon_i}
 \le\frac{\min\{R_j^2,B^2\}}{\sqrt n}
 \le\frac{BR_j}{\sqrt n}\le Bh_nR_j.
\]
Its centered mean is therefore at most
\(C\delta_zR_j\), and its centered range is at most \(2B^2\).

For each process and every \(0\le j\le J\), apply
Lemma~\ref{lem:supremum-bernstein} with
\(x_j=z+2j+\log8\).  The multiplier and square-process failures form the
geometric sum
\begin{equation}
 4\sum_{j=0}^Je^{-x_j}
 \le\frac12e^{-z}\sum_{j=0}^\infty e^{-2j}<e^{-z}.
 \tag{B.2c}
\end{equation}
The shellwise variance and range terms satisfy, respectively,
\[
 R_j\sqrt{x_j/n}
 \le CR_j\{\sqrt{z/n}+\sqrt{j/n}+n^{-1/2}\}
 \le C\delta_zR_j,
\]
and
\begin{equation}
 Cx_j/n\le C\{z/n+j/n+1/n\}\le C\delta_z^2.
 \tag{B.2d}
\end{equation}
For \(j\ge1\), \(R_j<2\norm f_2\); on the inner ball,
\(R_0=\delta_z\).  These inequalities prove (B.3)--(B.4).
Finally, \(\delta_z\ge r_0\), so the sub-\(r_0\) region is contained in
the actual positive-radius \(\delta_z\)-ball.  No function is rescaled, and
no star-shapedness or scaling closure is used.  The same calculation covers
\(n=1\), zero stochastic budgets, \(z=1\), the maximum-radius shell, and the
case in which the inner ball is the whole class.  Countability makes every
supremum measurable.
\end{proof}

Apply Lemma~\ref{lem:localized-process} before inserting fitted indices.  Its
complete primitive call list is:
\begin{enumerate}
\item \(\mathsf P_{\pi,2}^{u}(z)\): on \(\cD_2\), the countable
represented core of \(\Delta\cG_\pi\), with multiplier \(u\) and with the
square process;
\item \(\mathsf P_{\mu,2}^{\rm sq}(z)\): on \(\cD_2\), the countable
represented core of \(\Delta\cG_\mu\), retaining the square process;
\item \(\mathsf P_{J,2}^{\varepsilon}(z)\): on \(\cD_2\), the joint
class \(\mathcal J_\mu\) defined below, with multiplier \(\varepsilon\)
and with the square process;
\item \(\mathsf P_{\mu,1}^{u}(z)\): on \(\cD_1\), the countable
represented core of \(\Delta\cG_\mu\), with multiplier \(u\) and with the
square process;
\item \(\mathsf P_{\pi,2}^{\xi}(z)\): on \(\cD_2\), the countable
represented core of \(\Delta\cG_\pi\), with multiplier
\(\xi=Y-m_0(X)\);
\item \(\mathsf P_{\pi,2}^{m_0}(z)\): on \(\cD_2\), the countable
represented core of \(\Delta\cG_\pi\), with the fixed multiplier
\(m_0(X)\); and
\item \(\mathsf P_{\pi,1}^{m_0}(z)\): on \(\cD_1\), the countable
represented core of \(\Delta\cG_\pi\), with the fixed multiplier
\(m_0(X)\) and with the square process.
\end{enumerate}
The model-field certificate (A.1) is an outer-expectation bound for the full
difference class.  The pointwise inclusion in Appendix~\ref{app:model}
therefore gives the ordinary countable-core bound needed in (B.2a).  The
seven calls use the following explicit core base radii:
\begin{center}
\begin{tabular}{ll}
\(\mathsf P_{\pi,2}^{u}(z)\) & \(3t\) \\
\(\mathsf P_{\mu,2}^{\rm sq}(z)\) & \(3s\) \\
\(\mathsf P_{J,2}^{\varepsilon}(z)\) & \(C_J(s\vee h_n)\) \\
\(\mathsf P_{\mu,1}^{u}(z)\) & \(3s\) \\
\(\mathsf P_{\pi,2}^{\xi}(z)\) & \(3t\) \\
\(\mathsf P_{\pi,2}^{m_0}(z)\) & \(3t\) \\
\(\mathsf P_{\pi,1}^{m_0}(z)\) & \(3t\) \\
\end{tabular}
\end{center}
Indeed, (A.1) gives (B.2a) with \(r_0=3s\) or \(r_0=3t\)
for each simple core.  Consequently
\[
 K_1(3s\vee h_n)\le3K_1(s\vee h_n),\qquad
 K_1(3t\vee h_n)\le3K_1(t\vee h_n).
\]
Since \(C_J\ge4\), the already frozen \(C_{\rm rad}\) in the opening
display absorbs both \(3K_{1,\star}\) and \(K_{1,\star}C_J\); no downstream
tuning is changed.
If \(s=0\) or \(t=0\), the relevant call uses \(r_0=3s=0\) or
\(r_0=3t=0\) and the zero full-class certificate directly.  There is no radius-zero quotient.
Repeated downstream uses reuse these uniform events; they are not new calls
at fitted random indices.  Let
\[
 \mathcal F_2=\sigma(\cD_2),\quad
 \mathcal F_X=\mathcal F_2\vee\sigma(X_{1:n}),\quad
\mathcal F_T=\mathcal F_X\vee\sigma(T_{1:n}).
\]
These primitive calls are unconditional iid empirical-process calls.  The
nested sigma-fields are used later for the fitted weighted inequality
(B.1a) and, in the case of \(\mathcal F_2\), for the R denominator:
given \(\mathcal F_X\), the \(u_i\)'s remain independent and centered, and
given \(\mathcal F_T\), the \(\varepsilon_i\)'s do.  Let
\(\mathcal E_z^{\rm proc}\) be the finite intersection of the seven listed
uniform process events.  Apply every primitive call at the same admissible
confidence parameter \(z\); a finite union bound absorbs the seven fixed
failure multipliers into \(K_4\) and gives
\begin{equation}
 \Pp\{(\mathcal E_z^{\rm proc})^c\}\le K_4e^{-z},
 \qquad 1\le z\le n.
 \tag{B.5}
\end{equation}
This localized-process event is constructed separately for each \(z\);
simultaneity in \(z\) is not used.  It does not contain the comparator,
scalar-Bernstein, denominator, or fitted weighted events defined below.

\begin{lemma}[Sequential M1 transfer]
\label{lem:sequential-density-transfer}
Let \(\cG^{\rm rep}=\{g_j:j\in\mathbb N\}\subset\cG\) obey the common
envelope and suppose that, for every \(g\in\cG\), some represented sequence
converges pointwise everywhere to \(g\).  Then, for every \(P_X\), bounded
truth, and finite sample: (i) full-class and represented population infima
coincide; (ii) full-class and represented empirical infima coincide for every
continuous fit objective used here; (iii) full and represented
difference-class suprema coincide for the calibration objectives; and (iv)
the ordered first-within-tolerance represented selector is Borel.
\end{lemma}
\begin{proof}
The common envelope dominates squared differences, so pointwise convergence
and dominated convergence give \(L_2(P_X)\) convergence and hence (i).
Convergence of the finite evaluation vectors, followed by continuity of
finite sums, products, absolute values, and squares, gives (ii).  Pair
representing sequences for \(g,g'\) to approximate \(g-g'\); the same
finite-vector argument gives (iii).  Finally, if \(L_j(D)\) are the countable
Borel objectives, then \(L_*(D)=\inf_jL_j(D)\) is Borel and the least index
with \(L_j(D)\le L_*(D)+\tau_n\) has Borel level sets, proving (iv).
\end{proof}

For every empirical infimum over a public class, use its supplied countable
representation and select the first index within \(\tau_n\).  A countable
infimum of Borel functions and its first qualifying index are Borel.  The
set is nonempty by the definition of infimum, even without attainment.
Lemma~\ref{lem:sequential-density-transfer} transfers all full-class
population and finite-sample objectives used below to that representation.

\subsection{Initial fits}
\label{app:upper-initial}

On \(\cD_2\), let \(\widehat\pi\) be the first
\(\tau_n\)-minimizer of \(P_2(T-g)^2\) over \(\cG_\pi\), and let
\((\widehat\beta_0,\widehat\mu)\) be the first \(\tau_n\)-minimizer of
\(P_2(Y-\beta T-g)^2\) over the rational representation of
\([-3,3]\times\cG_\mu\).  For arbitrary
\(\varepsilon_{\rm cmp,\pi}>0\),
Lemma~\ref{lem:sequential-density-transfer} supplies a deterministic
represented comparator \(g_\pi\) satisfying
\(\norm{g_\pi-\pi_0}_2\le b+\varepsilon_{\rm cmp,\pi}\).  With
\(h=\widehat\pi-g_\pi\) and \(r_\pi=g_\pi-\pi_0\), the basic inequality is
\[
 P_2h^2\le2P_2(uh)-2P_2\{(g_\pi-\pi_0)h\}+\tau_n.
 \tag{B.6}
\]
For this fixed comparator, define the one-sided scalar-Bernstein event
\[
\begin{split}
 \mathcal E_{\pi,z}^{\rm cmp}(\varepsilon_{\rm cmp,\pi})
 :=\bigg\{P_2r_\pi^2\le{}&\norm{r_\pi}_2^2
 +B_\pi\norm{r_\pi}_2\sqrt{2z/n}
 +\frac{2B_\pi^2z}{3n}\bigg\}.
\end{split}
\]
Indeed, \(r_\pi^2\le B_\pi^2\) and
\(\operatorname{Var}(r_\pi^2)\le B_\pi^2\norm{r_\pi}_2^2\), so its
failure is at most \(2e^{-z}\), uniformly in the positive slack.  On
\(\mathcal E_z^{\rm proc}\cap
\mathcal E_{\pi,z}^{\rm cmp}(\varepsilon_{\rm cmp,\pi})\),
put \(w_z=\sqrt{z/n}\) and
\[
 \delta_{\pi,z}:=K_{1,\star}(3t\vee h_n)+w_z
 \le \bar t+w_z.
\]
Because \(h\) belongs to the represented treatment difference core and
\(P(uh)=0\), equations (B.3)--(B.4) give
\begin{equation}
 \begin{split}
 |P_2(uh)|&\le C_\rho\{\delta_{\pi,z}\norm h_2
                         +\delta_{\pi,z}^2\},\\
 P_2h^2&\ge \norm h_2^2
       -C_\rho\{\delta_{\pi,z}\norm h_2+\delta_{\pi,z}^2\}.
 \end{split}
 \tag{B.6a}
\end{equation}
Writing \(R_\pi=\norm{r_\pi}_{2,n}\), the fixed-comparator event and
\(\norm{r_\pi}_2\le b+\varepsilon_{\rm cmp,\pi}\) imply
\begin{equation}
 R_\pi\le C_\rho\{b+\varepsilon_{\rm cmp,\pi}+w_z\}.
 \tag{B.6b}
\end{equation}
Indeed, this follows by taking square roots in the displayed Bernstein bound
and using \(\sqrt{x^2+xy+y^2}\le x+y\) after absorbing the fixed envelope
coefficients into \(C_\rho\).  The upper side of (B.4) similarly gives the
population-to-empirical square-process conversion
\begin{equation}
 \norm h_{2,n}\le C_\rho(\norm h_2+\delta_{\pi,z}).
 \tag{B.6c}
\end{equation}
Thus (B.6), (B.6a), empirical Cauchy--Schwarz, and (B.6c) yield
\begin{align}
 \norm h_2^2
 &\le C_\rho\{\delta_{\pi,z}\norm h_2+\delta_{\pi,z}^2\}
      +2R_\pi\norm h_{2,n}+\tau_n \nonumber\\
 &\le C_\rho\{(\delta_{\pi,z}+R_\pi)\norm h_2
              +(\delta_{\pi,z}+R_\pi)^2+\tau_n\}.
 \tag{B.6d}
\end{align}
Specifically, Young's inequality gives
\[
 C_\rho(\delta_{\pi,z}+R_\pi)\norm h_2
 \le \frac12\norm h_2^2
      +C_\rho(\delta_{\pi,z}+R_\pi)^2,
\]
so the quadratic term is absorbed into the left-hand side.  Using
\(\sqrt{\tau_n}=\rho_n\), (B.6b), and
\(\delta_{\pi,z}\le\bar t+w_z\) gives
\begin{equation}
 \norm h_2\le C_\rho\{b+\varepsilon_{\rm cmp,\pi}
                         +\bar t+w_z+\rho_n\}.
 \tag{B.6e}
\end{equation}
Finally, the triangle inequality
\(\norm{\widehat\pi-\pi_0}_2\le\norm h_2+\norm{r_\pi}_2\)
gives, with a failure bound independent of
\(\varepsilon_{\rm cmp,\pi}\),
\[
 \norm{\widehat\pi-\pi_0}_2
 \le K_5\{b+t+\bar t+\sqrt{z/n}+\rho_n\}
      +C_\rho\varepsilon_{\rm cmp,\pi}.
\]
The resulting scalar tail inequality is uniform in the positive comparator
slack.  Letting \(\varepsilon_{\rm cmp,\pi}\downarrow0\) by continuity of
the scalar tail probabilities gives
\begin{equation}
 \norm{\widehat\pi-\pi_0}_2
 \le K_5\{b+t+\bar t+\sqrt{z/n}+\rho_n\}.
 \tag{B.7}
\end{equation}
with failure at most \(K_{\pi,2}e^{-z}\).  This limit does not intersect
the comparator-specific events.

For the joint fit first record the population curvature.  For every
\(\gamma\in\mathbb R\) and measurable \(f:\mathcal X\to\mathbb R\),
conditional centering gives
\[
 P(\gamma T+f)^2=\gamma^2\E u^2+\norm{\gamma\pi_0+f}_2^2
 \ge \frac{v_-}{2(1+B_0)^2}\{\gamma^2+\norm f_2^2\}.
\tag{B.8}
\]
Indeed, if \(q=\norm{\gamma\pi_0+f}_2\), then
\[
 \gamma^2+\norm f_2^2
 \le (1+2B_0^2)\gamma^2+2q^2.
\]
For \(c=v_-/\{2(1+B_0)^2\}\), one has
\(c(1+2B_0^2)\le v_-\le v\) and \(2c<1\), so
\(c\{\gamma^2+\norm f_2^2\}\le v\gamma^2+q^2\), proving (B.8).
The transfer to a joint class requires a separate localization argument.
Let \(\Delta\cG_\mu^{\rm rep}\) be the supplied countable difference core and
define
\[
 \mathcal J_\mu=\{\gamma T+f(X):
 \gamma\in\mathbb Q\cap[-6,6],\ f\in\Delta\cG_\mu^{\rm rep}\}.
\]
This class is countable, contains zero, and has the fixed envelope
\(B_{\mathcal J}=6B_0+2B_{\mathcal G}\).  If
\(j=\gamma T+f(X)\) and \(\norm j_2\le R\), then the orthogonal identity
\[
 \norm j_2^2=v\gamma^2+\norm{\gamma\pi_0+f}_2^2
\]
implies
\(\abs\gamma\le R/\sqrt{v_-}\) and \(\norm f_2\le C_fR\).  Hence, for
\(R\ge r_{0,\mu}:=C_J(s\vee h_n)\),
\[
\begin{aligned}
 \mathfrak R_{n,P}(R;\mathcal J_\mu)
 &\le \frac{R}{\sqrt{v_-}}\E\left|\frac{1}{n}
       \sum_{i=1}^n\epsilon_iT_i\right| \\
 &\quad +\mathfrak R_{n,P_X}(C_fR;\Delta\cG_\mu^{\rm rep}) \\
 &\le\left\{\frac{B_0}{\sqrt{v_-}}h_n
             +3c_{\rm rad}C_fs\right\}R \\
 &\le c_{\rm rad}r_{0,\mu}R,
\end{aligned}
\]
where the last inequality follows from the definition of \(C_J\).  This uses
the public outcome-class budget only at admissible positive radii; when
\(s=0\), the floor \(C_Jh_n\) remains positive.
Lemma~\ref{lem:localized-process}, with
base observation \((X,T)\) and conditionally centered multiplier
\(\varepsilon\), therefore supplies both the multiplier and square-process
bounds uniformly over \(\mathcal J_\mu\).

For completeness, here is the full basic-inequality step.  For arbitrary
\(\varepsilon_{\rm cmp,\mu}>0\),
Lemma~\ref{lem:sequential-density-transfer} supplies deterministic represented
comparators \(\beta^\circ\in\mathbb Q\cap[-3,3]\) and
\(g^\circ\in\cG_\mu\) satisfying
\(\abs{\beta^\circ-\beta_0}\le\varepsilon_{\rm cmp,\mu}\) and
\(\norm{g^\circ-\mu_0}_2\le a+\varepsilon_{\rm cmp,\mu}\), and put
\[
 e^\circ=(\beta^\circ-\beta_0)T+(g^\circ-\mu_0),\qquad
 j=(\widehat\beta_0-\beta^\circ)T+(\widehat\mu-g^\circ).
\]
The approximate empirical minimization gives exactly
\[
 P_2j^2\le2P_2(\varepsilon j)-2P_2(e^\circ j)+\tau_n.
\]
For this fixed pair of comparators, let
\(B_{\rm cmp,\mu}:=7B_0+B_{\mathcal G}\) and define
\[
\begin{split}
 \mathcal E_{\mu,z}^{\rm cmp}(\varepsilon_{\rm cmp,\mu})
 :=\bigg\{P_2(e^\circ)^2\le{}&\norm{e^\circ}_2^2
 +B_{\rm cmp,\mu}\norm{e^\circ}_2\sqrt{2z/n}
 +\frac{2B_{\rm cmp,\mu}^2z}{3n}\bigg\}.
\end{split}
\]
Here \(\abs{e^\circ}\le B_{\rm cmp,\mu}\) and
\(\operatorname{Var}\{(e^\circ)^2\}\le
B_{\rm cmp,\mu}^2\norm{e^\circ}_2^2\), so its failure is at most
\(2e^{-z}\), uniformly in the positive slack.
Put \(w_z=\sqrt{z/n}\) and
\(\delta_{\mu,z}=C_{\rm rad}(s\vee h_n)+w_z\).  Since \(j\) belongs to the
countable joint class \(\mathcal J_\mu\), the multiplier and square parts of
\(\mathsf P_{J,2}^{\varepsilon}(z)\), before insertion of the fitted index,
give, using \(P(\varepsilon j)=0\),
\begin{equation}
 \begin{split}
 |P_2(\varepsilon j)|+|(P_2-P)j^2|
 &\le C_\rho\{\delta_{\mu,z}\norm j_2+\delta_{\mu,z}^2\},\\
 \norm j_{2,n}&\le C_\rho(\norm j_2+\delta_{\mu,z}).
 \end{split}
 \tag{B.8a}
\end{equation}
The deterministic comparator residual obeys
\[
 \norm{e^\circ}_2
 \le B_0|\beta^\circ-\beta_0|+\norm{g^\circ-\mu_0}_2
 \le a+(1+B_0)\varepsilon_{\rm cmp,\mu}.
\]
Consequently the fixed-comparator Bernstein event gives, for
\(E_\mu:=\norm{e^\circ}_{2,n}\),
\begin{equation}
 E_\mu\le C_\rho\{a+(1+B_0)\varepsilon_{\rm cmp,\mu}+w_z\}.
 \tag{B.8b}
\end{equation}
The joint basic inequality, the lower square-process comparison in (B.8a),
and empirical Cauchy--Schwarz now give
\begin{align}
 \norm j_2^2
 &\le |(P_2-P)j^2|+2|P_2(\varepsilon j)|
      +2E_\mu\norm j_{2,n}+\tau_n \nonumber\\
 &\le C_\rho\{(\delta_{\mu,z}+E_\mu)\norm j_2
              +(\delta_{\mu,z}+E_\mu)^2+\tau_n\}.
 \tag{B.8c}
\end{align}
Here the explicit quadratic absorption is
\[
 C_\rho(\delta_{\mu,z}+E_\mu)\norm j_2
 \le\frac12\norm j_2^2
     +C_\rho(\delta_{\mu,z}+E_\mu)^2.
\]
Together with \(\sqrt{\tau_n}=\rho_n\), it yields
\begin{equation}
 \norm j_2\le C_\rho\{a+(1+B_0)\varepsilon_{\rm cmp,\mu}
                         +\bar s+w_z+\rho_n\}.
 \tag{B.8d}
\end{equation}
Set \(\Gamma=\widehat\beta_0-\beta_0\) and
\(F=\widehat\mu-\mu_0\).  Since \(\Gamma T+F=e^\circ+j\), the population
comparator bound, the triangle inequality, and (B.8) imply
\begin{equation}
 \begin{split}
 |\Gamma|+\norm F_2
 &\le \sqrt2\{\Gamma^2+\norm F_2^2\}^{1/2}\\
 &\le \frac{2(1+B_0)}{\sqrt{v_-}}\norm{\Gamma T+F}_2\\
 &\le C_\rho\{a+(1+B_0)\varepsilon_{\rm cmp,\mu}
                +\bar s+w_z+\rho_n\}.
 \end{split}
 \tag{B.8e}
\end{equation}
Equations (B.8a)--(B.8e) use only the already listed joint process call and
the fixed-comparator event; they add no joint-complexity certificate.  Their
failure bound is independent of the positive comparator slack.  Letting
\(\varepsilon_{\rm cmp,\mu}\downarrow0\) by continuity of the scalar tail
probabilities
yields
\begin{equation}
 \abs{\widehat\beta_0-\beta_0}+\norm{\widehat\mu-\mu_0}_2
 \le K_6\{a+s+\bar s+\sqrt{z/n}+\rho_n\}.
 \tag{B.9}
\end{equation}
with failure at most \(K_{\mu,2}e^{-z}\).  Again, no infinite
comparator-event intersection is formed.
The proof-only comparator slack is not an estimator tuning parameter.  Thus
the proof uses neither an uncountable fitted-index insertion nor a new joint
complexity assumption.
No exact minimizer, convexity, or population oracle is used by the learner.

\paragraph{Empirical initial-fit radii.}
Let \(K_{\rm sq}\) be the maximum, after the preceding radius rescaling, of
the square-process coefficients for the represented treatment and outcome
difference cores, and put
\[
 c_{\rm sq}=K_{\rm sq}+\sqrt{K_{\rm sq}},\qquad
 L_\pi=K_5+c_{\rm sq}+B_\pi+2,\qquad
 L_\mu=K_6+c_{\rm sq}+B_\pi+2.
\]
The square-process event is uniform over each countable difference core before
the fitted index is inserted and implies, on either named half,
\(\norm{f}_{j,n}\le\norm f_2+c_{\rm sq}\Delta\).  For a deterministic
comparator residual \(r\), bounded by \(B_\pi\), scalar Bernstein gives
\(\norm{r}_{j,n}\le\norm r_2+C_\rho B_\pi\sqrt{z/n}\).
Applying this first on the fitting half \(\cD_2\), for every positive
comparator slack,
\[
 \|\widehat\pi-\pi_0\|_{2,n}
 \le L_\pi\{b+t+\bar t+\sqrt{z/n}+\rho_n\}
       +2\varepsilon_{\rm cmp,\pi},
\]
and
\[
 |\widehat\beta_0-\beta_0|+\|\widehat\mu-\mu_0\|_{2,n}
 \le L_\mu\{a+s+\bar s+\sqrt{z/n}+\rho_n\}
       +2\varepsilon_{\rm cmp,\mu}.
\]
For each display the failure bound is independent of the slack.  Taking the
corresponding slack down to zero in the scalar tail probabilities, rather than intersecting comparator-specific events, yields the same displays without the
final slack terms.  Define \(\mathcal E_{\pi,z}^{\rm fit,2}\) to be the event
on which (B.7) and the exact \(\cD_2\) empirical treatment-fit display both
hold, and define \(\mathcal E_{\mu,z}^{\rm fit,2}\) analogously using (B.9).
Their failure bounds are
\[
 \Pp\{(\mathcal E_{\pi,z}^{\rm fit,2})^c\}
 \le K_{\pi,\rm fit,2}e^{-z},\qquad
 \Pp\{(\mathcal E_{\mu,z}^{\rm fit,2})^c\}
 \le K_{\mu,\rm fit,2}e^{-z}.
\]
These are events defined by the displayed exact inequalities, not
intersections over comparator choices.

The downstream edits instead require \(\cD_1\) empirical norms.  Condition on
\(\mathcal F_2=\sigma(\cD_2)\).  Then each fitted represented index and hence
\(h=\widehat\pi-g_\pi\) and \(f=\widehat\mu-g^\circ\) is conditionally fixed,
while \(\cD_1\) remains iid.  Before conditioning, the uniform square parts of
\(\mathsf P_{\pi,1}^{m_0}(z)\) and
\(\mathsf P_{\mu,1}^{u}(z)\) were already formed over their entire represented
difference cores.  Thus they may be evaluated at these conditionally fixed
indices and give
\[
 \norm{h}_{1,n}\le\norm h_2+c_{\rm sq}
   \{C_{\rm rad}(t\vee h_n)+\sqrt{z/n}\},\qquad
 \norm{f}_{1,n}\le\norm f_2+c_{\rm sq}
   \{C_{\rm rad}(s\vee h_n)+\sqrt{z/n}\}.
\]
For each fixed positive proof slack, apply scalar Bernstein on \(\cD_1\) to
the fixed comparator residuals \(r_\pi=g_\pi-\pi_0\) and
\(r_\mu=g^\circ-\mu_0\).  Denote the resulting events by
\(\mathcal E_{\pi,z}^{\rm cmp,1}(\varepsilon_{\rm cmp,\pi})\) and
\(\mathcal E_{\mu,z}^{\rm cmp,1}(\varepsilon_{\rm cmp,\mu})\); each has
failure at most \(2e^{-z}\), \(1\le z\le n\), and supplies the preceding
fixed-residual norm conversion.  Combining these events with (B.7)--(B.9),
the triangle inequality, and then taking each proof slack down to zero by
scalar tail-probability continuity gives
\begin{align}
 \norm{\widehat\pi-\pi_0}_{1,n}
 &\le L_\pi\{b+t+\bar t+\sqrt{z/n}+\rho_n\},\tag{B.9a}\\
 |\widehat\beta_0-\beta_0|+
 \norm{\widehat\mu-\mu_0}_{1,n}
 &\le L_\mu\{a+s+\bar s+\sqrt{z/n}+\rho_n\}.
 \tag{B.9b}
\end{align}
No infinite comparator-event intersection is used, and neither slack is an
estimator input.  Let \(\mathcal E_{\pi,z}^{\rm tr,1}\) and
\(\mathcal E_{\mu,z}^{\rm tr,1}\) be the exact events defined respectively
by (B.7),(B.9a) and by (B.9),(B.9b).  The finite unions just described give
\[
 \Pp\{(\mathcal E_{\pi,z}^{\rm tr,1})^c\}
 \le K_{\pi,\rm tr,1}e^{-z},\qquad
 \Pp\{(\mathcal E_{\mu,z}^{\rm tr,1})^c\}
 \le K_{\mu,\rm tr,1}e^{-z},\qquad 1\le z\le n.
\]
These are population-to-\(\cD_1\) transfer events; they are distinct from
the \(\cD_2\) fit events above and change only fixed failure multipliers.

\subsection{The O edit: definition, normalization, and calibration}
\label{app:upper-o-calibration}

Put \(x=a+s\), \(y=b+t\), and recall the frozen envelope \(B_\pi\).  Put
\[
 d_\pi=\min\{B_\pi,8(y+\bar s+\rho_n)\},\qquad
 \lambda_O=(\bar s/d_\pi)^2.
\]
The cap is fixed because
\(\abs{\widehat\pi-\pi_0}\le B_{\mathcal G}+B_0=B_\pi\)
pointwise.
On the O-safe cell
\[
 x\le y,\qquad \bar s\le\vartheta_\rho,
\]
both arguments in the minimum are at least \(8\bar s\).  Hence
\[
 d_\pi\ge8\bar s,\qquad0<\lambda_O\le1/64<1/16.
\]
Define
\[
 \widetilde D_O=(v_-/2)\vee P_1\{T(T-\widehat\pi)\},\qquad
 \widetilde a_i=\frac{T_i-\widehat\pi(X_i)}{\widetilde D_O}.
\]
For \(f\in\Delta\cG_\mu\), \(\gamma\in\R\), and \(a\in\R^n\), let
\begin{align*}
 \Phi_\mu(f,\gamma,a)
 &=\abs{P_1\{(\gamma T+f)a\}-\gamma}-P_1f^2,\\
 H_\mu(a)&=\sup_{\gamma\in\R,\,f\in\Delta\cG_\mu}
 \Phi_\mu(f,\gamma,a),\\
 F_O(a)&=\lambda_O\norm{a-\widetilde a}_n^2+H_\mu(a).
\end{align*}
Here and below, the code takes the supremum over the countable represented
core.  On a finite sample, Lemma~\ref{lem:sequential-density-transfer}(iii)
identifies its value with the displayed full-class supremum.
Because \(0\in\Delta\cG_\mu\), \(H_\mu(a)<\infty\) if and only if
\(P_1(Ta)=1\).  On this hyperplane,
\[
 H_\mu(a)=\sup_f\{\abs{P_1(fa)}-P_1f^2\}.
 \tag{B.10}
\]
For \(a,a'\) on this normalization hyperplane, the uniform envelope gives
\[
 |H_\mu(a)-H_\mu(a')|
 \le\sup_{f\in\Delta\cG_\mu}|P_1\{f(a-a')\}|
 \le2B_{\mathcal G}\norm{a-a'}_n.
\]
Hence the reduced objective is continuous on the hyperplane; no continuity
of the extended-valued \(H_\mu\) on all of \(\R^n\) is asserted.
If \(T_{1:n}=0\), set
\[
 \widehat a:=0,\qquad
 \widehat\beta_O^{\rm raw}:=0.
\]
On \(T_{1:n}\ne0\), choose the first largest-magnitude coordinate \(j\),
enumerate rational values of the other coordinates, and set
\[
 a_j=\frac{n-\sum_{i\ne j}T_ia_i}{T_j}.
 \tag{B.11}
\]
This is a Borel dense chart of the exact hyperplane, even for irrational
\(T_j\).  The objective is continuous and coercive there.  The first point
within \(\tau_n\) of its countable infimum is a Borel
\(\tau_n\)-minimizer \(\widehat a\) satisfying \(P_1(T\widehat a)=1\), and
we put
\[
 \widehat\beta_O^{\rm raw}=P_1\{(Y-\widehat\mu)\widehat a\}.
\]
Finally define
\[
 \widehat\beta_O=\Pi_{[-3,3]}(\widehat\beta_O^{\rm raw}).
\]
With \(p_\rho=v_-/B_0^2\),
\[
 \E T^2=\E\pi_0(X)^2+v\ge v_-,\qquad
 \Pr(T\ne0)\ge\frac{\E T^2}{B_0^2}\ge p_\rho.
\]
Independence of the D1 observations therefore gives
\[
 \Pr(T_{1:n}=0)\le e^{-p_\rho n},
\]
and projection bounds the loss on this event by three.  Therefore
\[
 3\Pr(T_{1:n}=0)
 \le 3e^{-p_\rho n}
 \le \frac{3}{\sqrt{2ep_\rho}}\,n^{-1/2}.
\]
The O proof below is on \(T_{1:n}\ne0\); this display accounts separately
for the zero-vector branch.

For the proof-only comparator let \(D_0=P_1(Tu)\) and \(a_0=u/D_0\).
Since \(\E(Tu)=v\), \(\abs{Tu}\le B_0^2\), and
\(\operatorname{Var}(Tu)\le \E(T^2u^2)\le B_0^2v\), define
\[
 \mathcal E_{D_0,z}:=
 \left\{\abs{D_0-v}\le
 B_0\sqrt{\frac{2vz}{n}}+\frac{4B_0^2z}{3n}\right\}.
\]
Lemma~\ref{lem:bernstein} gives
\(\Pp(\mathcal E_{D_0,z}^c)\le2e^{-z}\): the centered range satisfies
\(\abs{Tu-v}\le2B_0^2\), which gives the displayed \(4/3\) linear
coefficient.  With
\(c_0=v_-/(32B_0^2)\), the square-root and linear deviations are at most
\(v/4\) and \(v/24\), respectively, whenever
\(1\le z\le c_0n\).  Hence \(D_0\ge17v/24\ge v_-/2\) on
\(\mathcal E_{D_0,z}\), and \(P_1(Ta_0)=1\).  On
\(\mathcal E_z^{\rm proc}\cap\mathcal E_{D_0,z}\), apply the localized
multiplier and square bounds to (B.10).  Because
\(\E\{f(X)u\}=0\), with \(w_z=\sqrt{z/n}\) they give
\[
 \abs{P_1(fa_0)}\le C_\rho\{(\bar s+w_z)\norm f_2
                  +(\bar s+w_z)^2\},
 \quad
 P_1f^2\ge\norm f_2^2-C_\rho\{(\bar s+w_z)\norm f_2
                  +(\bar s+w_z)^2\}.
\]
Maximizing the resulting concave quadratic gives the all-\(z\) benchmark
\begin{equation}
 H_\mu(a_0)\le K_7\{\bar s^2+z/n\}.
 \tag{B.12}
\end{equation}
Projection of the fitted denominator onto \([v_-/2,\infty)\) is
one-Lipschitz and leaves \(D_0\) fixed.  With
\(e_\pi=\widehat\pi-\pi_0\),
\[
 \abs{\widetilde D_O-D_0}\le\abs{P_1(Te_\pi)}
 \le B_0\norm{e_\pi}_{1,n}.
\]
On \(D_0\ge v_-/2\), put
\[
 C_u:=\frac2{v_-}+\frac{4B_0^2}{v_-^2}.
\]
The same denominator algebra gives
\[
 \|\widetilde a-a_0\|_n
 \le C_u\|\widehat\pi-\pi_0\|_{1,n}.
\]
If the cap in \(d_\pi\) is inactive,
 \(\mathcal E_{\pi,z}^{\rm tr,1}\) and (B.9a),
\(\bar t\le C_{\rm rad}y+\bar s\) bound the deterministic part by a fixed
multiple of \(d_\pi\).  If the cap is active, the pointwise envelope gives
\(\abs{\widehat\pi-\pi_0}\le B_\pi=d_\pi\).  Put
\[
 K_8^*:=C_u\max\left\{1,L_\pi,
              \frac{L_\pi(1+C_{\rm rad})}{8}\right\}.
\]
Thus, including cap equality,
\begin{equation}
 \norm{\widetilde a-a_0}_n
 \le K_8^*\{d_\pi+\sqrt{z/n}\}.
 \tag{B.13}
\end{equation}
This covers active and inactive caps.

Compare the approximate objective at \(\widehat a\) and \(a_0\).  Since
\(H_\mu\ge0\),
\begin{align}
 \norm{\widehat a-\widetilde a}_n
 &\le K_9\{d_\pi+\sqrt{z/n}
      +(d_\pi/\bar s)\sqrt{z/n}
      +\sqrt{\tau_n/\lambda_O}\},\tag{B.14}\\
 H_\mu(\widehat a)
 &\le K_9\{\bar s^2+z/n
       +\lambda_O(d_\pi+\sqrt{z/n})^2+\tau_n\}.
 \tag{B.15}
\end{align}
Indeed, approximate optimality first bounds
\(\lambda_O\norm{\widehat a-\widetilde a}_n^2+H_\mu(\widehat a)\) by
\(K_\rho\{\lambda_O(d_\pi+w_z)^2+\bar s^2+w_z^2+\tau_n\}\);
only the norm bound divides by \(\lambda_O\), and
\(\sqrt{\lambda_O}\asymp_\rho\bar s/d_\pi\).  This is the source of the
amplified term in (B.14).
Exact normalization and (B.10) now give, for every represented \(f\),
\begin{equation}
 \abs{P_1(f\widehat a)}
 \le P_1f^2+H_\mu(\widehat a).
 \tag{B.16}
\end{equation}
Thus the one-sided O calibration inequality has been proved rather than
assumed.
Equations (B.12)--(B.16) hold on the finite intersection
\(\mathcal E_z^{\rm proc}\cap\mathcal E_{D_0,z}\cap
\mathcal E_{\pi,z}^{\rm tr,1}\), for
\(1\le z\le(c_0\wedge1)n\); each fitted substitution is made only after
the corresponding uniform or exact-fit event has been formed.

\subsection{The O all-z oracle inequality and expectation}
\label{app:upper-o-oracle}

Recall the piecewise raw O output and its clipped version defined above.
Choose a deterministic represented \(g_\mu\), available by
Lemma~\ref{lem:sequential-density-transfer}, such that
\(\norm{g_\mu-\mu_0}_2\le a+\rho_n\), and put \(w=u/v\) and
\(r_\mu=\mu_0-g_\mu\).
Exact normalization gives the exhaustive identity
\begin{align}
 \widehat\beta_O^{\rm raw}-\beta_0
={}&P_1(\varepsilon w)+P_1\{\varepsilon(\widehat a-w)\}
 +P_1\{(\mu_0-g_\mu)w\}\nonumber\\
 &+P_1\{(\mu_0-g_\mu)(\widehat a-w)\}
 +P_1\{(g_\mu-\widehat\mu)\widehat a\}.
 \tag{B.17}
\end{align}
Let \(r=\bar s\), \(d=d_\pi\), \(w_z=\sqrt{z/n}\), and
\(\ell=\sqrt{\tau_n/\lambda_O}\).  On \(\mathcal E_{D_0,z}\), the
preceding deviation bound and \(\norm u_n\le B_0\) give
\(\norm{a_0-w}_n\le K_\rho w_z\), hence (B.13)--(B.14) imply
\[
 \norm{\widehat a-w}_n\le K\{d+w_z+(d/r)w_z+\ell\}.
\]
For this fixed comparator, let
\[
 \mathcal E_{O,z}^{\rm cmp}
 :=\{\norm{r_\mu}_{1,n}\le K_{O,\rm cmp}(a+\rho_n+w_z)\}.
\]
Bounded scalar Bernstein applied to \(r_\mu(X)^2\) gives
\(\Pp\{(\mathcal E_{O,z}^{\rm cmp})^c\}\le2e^{-z}\), after one fixed
enlargement of \(K_{O,\rm cmp}\).  No fitted index occurs in this event.
Define the fitted weighted event
\[
\begin{split}
 \mathcal E_{O,z}^{\rm wt}:=\bigg\{&
 \abs{P_1(\varepsilon w)}\le K_{O,\rm wt}w_z,\\
 &\abs{P_1\{\varepsilon(\widehat a-w)\}}
 \le K_{O,\rm wt}\norm{\widehat a-w}_n w_z,\\
 &\abs{P_1(r_\mu w)}
 \le K_{O,\rm wt}\norm{r_\mu}_{1,n} w_z\bigg\}.
\end{split}
\]
The first two inequalities follow from (B.1a) conditional on
\(\mathcal F_T\); the last follows from (B.1a) conditional on
\(\mathcal F_X\), because \(\E(u\mid X)=0\).  Taking expectations of
each conditional failure over its own conditioning field and then applying
an unconditional union bound gives failure at most \(6e^{-z}\).

On the finite event
\[
 \mathcal E_{O,z}:=
 \mathcal E_z^{\rm proc}\cap\mathcal E_{\pi,z}^{\rm tr,1}\cap
 \mathcal E_{\mu,z}^{\rm tr,1}\cap\mathcal E_{D_0,z}\cap
 \mathcal E_{O,z}^{\rm cmp}\cap\mathcal E_{O,z}^{\rm wt},
 \qquad \Pp(\mathcal E_{O,z}^c)\le K_{O,0}e^{-z},
\]
Equation (B.1a), the exact fit events, and (B.15)--(B.16) give the five
separate bounds
\begin{align*}
 \abs{P_1(\varepsilon w)}&\le Kw_z,\\
 \abs{P_1\{\varepsilon(\widehat a-w)\}}
   &\le K\{d+w_z+(d/r)w_z+\ell\}w_z,\\
 \abs{P_1\{(\mu_0-g_\mu)w\}}
   &\le K(aw_z+w_z^2),\\
 \abs{P_1\{(\mu_0-g_\mu)(\widehat a-w)\}}
   &\le K(a+w_z)\{d+w_z+(d/r)w_z+\ell\},\\
 \abs{P_1\{(g_\mu-\widehat\mu)\widehat a\}}
   &\le K\{(x+r+w_z)^2+r^2+w_z^2
             +\lambda_O(d+w_z)^2+\tau_n\}.
\end{align*}
Adding these inequalities gives the displayed bound first for the raw output.
Projection cannot increase distance to the legal target
\(\beta_0\in[-3,3]\), so the same tail bound holds for the clipped
\(\widehat\beta_O\).  Retaining a common nonnegative amplification envelope
for the subsequent layer-cake calculation gives
\begin{align}
 \abs{\widehat\beta_O-\beta_0}
 \le K_{10}\bigg[&ad+x^2+r^2+\lambda_Od^2+\tau_n+\ell(a+r)\nonumber\\
 &+(1+a+d+x+r)w_z+\frac{ad}{r}w_z
  +\left(1+\frac d r\right)w_z^2+w_z^3+\ell w_z\bigg]
 \tag{B.18}
\end{align}
on \(\mathcal E_{O,z}\).  Choose the fixed endpoint
\(c_2\le c_0\wedge1\) and \(K_{11}\ge K_{O,0}\); then its failure is at
most \(K_{11}e^{-z}\), \(1\le z\le c_2n\).
The displayed \(d w_z\) is the edited-noise contribution.  The amplified
terms \((ad/r)w_z\) and \((d/r)w_z^2\) are required by the corrected
intermediate norm bound (B.14), rather than decorative padding: respectively,
they arise when \((d/r)w_z\) is multiplied by \(a\) and by \(w_z\) in the
five-term identity.  There is no standalone \((d/r)w_z\), no
\((d/r)w_z^3\), and no \(\lambda_O^{-1}w_z^2\).  The denominator perturbation
is in (B.13) and is not counted twice.

\begin{lemma}[O all-\(z\) absorption certificate]
\label{lem:o-all-z-absorption}
On the O-safe cell, every monomial in (B.18), integrated against its
pointwise tail, is bounded by \(K(q+ay+x^2)\).  The complete accounting is
\[
\begin{array}{c|c|c}
\text{tail monomial}&\text{layer-cake scale}&\text{deterministic control}\\ \hline
(1+a+d+x+r)w_z&q&Kq\\
\frac{ad}{r}w_z&\frac{ad}{r}q&K(q+ay+x^2)\\
w_z^2&q^2&Kq\\
\frac d r w_z^2&\frac d{nr}&Kq\\
w_z^3&q^3&Kq\\
\ell w_z&\ell q&K(q+ay+x^2).
\end{array}
\tag{B.19}
\]
The deterministic terms \(ad,x^2,r^2,\lambda_Od^2,\tau_n\), and
\(\ell(a+r)\) obey the same target bound.  This statement includes zero
budgets, both cap branches, and cap equality.  It is pointwise in \(z\), not
simultaneous in \(z\); unsafe cells are handled only by the later fallback.
\end{lemma}
\begin{proof}
Use
\[
 h_n^2\le2q,\quad r\ge C_{\rm rad}h_n,\quad
 d\le B_\pi,\quad d\le K_d(y+r+\rho_n),\quad
 \lambda_Od^2=K_\lambda r^2.
\]
Because \(s\le x\le y\), \(r^2\le K(x^2+q)\),
\(\rho_n\le q\le Kr\), and \(ad\le K(ay+x^2+q)\): when
\(h_n\le y\), send \(ah_n\) to \(ay\); when \(y<h_n\), use
\(a\le x\le y<h_n\) and \(ah_n\le h_n^2\le2q\).
Since \(q/r\le K\), the same split proves
\((ad/r)q\le K(q+ay+x^2)\), while
\[
 \frac d{nr}\le
 \frac{B_\pi}{C_{\rm rad}\sqrt{n\log(en)}}\le Kq,\qquad
 \ell=\frac{\rho_n}{\sqrt{\lambda_O}}
 =\frac{\rho_nd}{r}.
\]
Thus \(\ell(a+r)+\ell q\) has the claimed bound.  The remaining rows use
\(q^2,q^3\le q\) and the fixed envelopes.  The positive floor
\(r=C_{\rm rad}(s\vee h_n)\) makes every ratio defined when \(s=0\), and
the two inequalities for \(d\) cover the active, inactive, and equality
cases of the cap.  Finally,
\(\int_1^\infty z^{p/2}e^{-z}\,dz\le\Gamma(1+p/2)\) gives the listed
scales.  The interval below \(z=1\) uses the \(z=1\) threshold; clipping
handles \(z>c_2n\); and if \(c_2n<1\), then \(q>\sqrt{c_2}\) absorbs the
bounded loss.  This uses a separate event for each \(z\), never one
simultaneous event.
\end{proof}

Applying Lemma~\ref{lem:o-all-z-absorption} to (B.18) and adding the zero-vector contribution
\(3/\sqrt{2ep_\rho}\,n^{-1/2}\) to the fixed O constant, we obtain
\begin{equation}
 \sup_{P\in\cP(\Gamma)}
 \E_P\abs{\widehat\beta_O-\beta(P)}
 \le C_O\{q+ay+x^2\}.
 \tag{B.20}
\end{equation}

\subsection{The R edit and two-sample calibration}
\label{app:upper-r-calibration}

This R proof is an independent replay, not by symmetry with O.  It uses only
the generic public-class process lemma and the corrected initial-fit bounds;
the O instrument, normalization, and functional \(H_\mu\) do not enter.
Let \(m_0(X)=\E(Y\mid X)=\beta_0\pi_0(X)+\mu_0(X)\),
and recall the frozen envelope \(B_{\rm all}\).  Put
\[
 d_{\rm all}=\min\{B_{\rm all},8(x+y+\bar t+\rho_n)\},\qquad
 \lambda_R=(\bar t/d_{\rm all})^2.
\]
Here \(\abs{m_0}\le B_0\) by conditional expectation of the bounded
response, while
\(\abs{\widehat\beta_0\widehat\pi+\widehat\mu}\le4B_{\mathcal G}\).
Thus \(B_{\rm all}\) is a fixed envelope for their discrepancy.
On the R-safe cell \(y<x\), \(\bar t\le\vartheta_\rho\), both arguments in
the minimum are at least \(8\bar t\), so
\[
 d_{\rm all}\ge8\bar t,\qquad0<\lambda_R\le1/64<1/16.
\]
On D1 covariates set
\(\widetilde m_i=\widehat\beta_0\widehat\pi(X_i)+\widehat\mu(X_i)\).
For \(f\in\Delta\cG_\pi\), define
\[
 H_\pi(m)=\sup_f\{\abs{P_2(Yf)-P_1(mf)}-P_1f^2\}.
\]
For every \(m,m'\in\R^n\), the uniform envelope similarly gives
\[
 |H_\pi(m)-H_\pi(m')|
 \le\sup_{f\in\Delta\cG_\pi}|P_1\{f(m-m')\}|
 \le2B_{\mathcal G}\norm{m-m'}_n.
\]
Thus \(H_\pi\) is continuous on \(\R^n\); together with the positive
quadratic penalty and coercivity below, this supports first-within-tolerance
selection on the rational dense chart.
The implemented supremum is over the represented difference core and equals
the full-class display by Lemma~\ref{lem:sequential-density-transfer}(iii).
Over the rational enumeration of \(\R^n\), select the first
\(\tau_n\)-minimizer \(\widehat m\) of
\(\lambda_R\norm{m-\widetilde m}_n^2+H_\pi(m)\).  The positive quadratic is
coercive, so this Borel selection is well defined.

Put \(\xi=Y-m_0(X)\) and
\(\delta_z=K(\bar t+\sqrt{z/n})\).  On
\(\mathcal E_z^{\rm proc}\), uniformly in
\(f\in\Delta\cG_\pi\), (B.3)--(B.4) give
\begin{align}
 &\abs{P_2(\xi f)}+\abs{(P_2-P)(m_0f)}
 +\abs{(P_1-P)(m_0f)}+\abs{(P_1-P)f^2}\nonumber\\
 &\hspace{35mm}\le K\{\delta_z\norm f_2+\delta_z^2\}.
 \tag{B.21}
\end{align}
Inserting \(m=m_0(X_{1:n})\) and maximizing
\(K\delta_zr-r^2/2\) proves
\begin{equation}
 H_\pi(m_0(X_{1:n}))\le K\delta_z^2.
 \tag{B.22}
\end{equation}
The exact decomposition
\[
 \widetilde m-m_0=(\widehat\beta_0-\beta_0)\widehat\pi
 +\beta_0(\widehat\pi-\pi_0)+(\widehat\mu-\mu_0)
\]
and \(\mathcal E_{\pi,z}^{\rm tr,1}\cap
\mathcal E_{\mu,z}^{\rm tr,1}\), together with (B.9a)--(B.9b) and
\(\bar s\le C_{\rm rad}x+\bar t\) on \(y<x\), give the following bound after
putting
\[
 L_R:=(B_{\mathcal G}+1)L_\mu(1+C_{\rm rad})+3L_\pi:
\]
\[
 \|\widetilde m-m_0\|_{1,n}
 \le L_R\{x+y+\bar t+\rho_n+\sqrt{z/n}\}.
\]
If the cap is inactive, the deterministic base is a fixed multiple of
\(d_{\rm all}\); if active, the pointwise envelope gives
\(\abs{\widetilde m-m_0}\le B_{\rm all}=d_{\rm all}\).  Put
\(K_m^*:=1+L_R\).  Hence, including equality,
\[
 \|\widetilde m-m_0\|_{1,n}\le K_m^*\{d_{\rm all}+\sqrt{z/n}\}.
\]
Objective comparison therefore yields
\begin{align}
 H_\pi(\widehat m)&\le K\{\delta_z^2+
 \lambda_R(d_{\rm all}+\delta_z)^2+\tau_n\},\tag{B.23}\\
 \norm{\widehat m-m_0}_n&\le K\{d_{\rm all}+\delta_z+
 \delta_z/\sqrt{\lambda_R}+\sqrt{\tau_n/\lambda_R}\}.
 \tag{B.24}
\end{align}
For every represented \(f\), the identity
\[
 P_1\{(m_0-\widehat m)f\}
 =[P_2(Yf)-P_1(\widehat mf)]-[P_2(Yf)-P_1(m_0f)]
\]
and (B.21)--(B.23) imply
\begin{equation}
 \abs{P_1\{(m_0-\widehat m)f\}}
 \le K\{\norm f_2^2+\delta_z^2+
 \lambda_R(d_{\rm all}+\delta_z)^2+\tau_n\}.
 \tag{B.25}
\end{equation}
The event is uniform before substitution, so (B.25) applies legally to
\(f=\widehat\pi-g_\pi\).  The edit uses D2 outcomes and D1 covariates only;
conditioning has not frozen the D1 treatments.
Equations (B.21)--(B.25) hold on
\(\mathcal E_z^{\rm proc}\cap\mathcal E_{\pi,z}^{\rm tr,1}\cap
\mathcal E_{\mu,z}^{\rm tr,1}\).

\subsection{The R denominator, oracle inequality, and expectation}
\label{app:upper-r-oracle}

Let \(e_\pi=\widehat\pi-\pi_0\), \(r_m=\widehat m-m_0\), and
\(D=P_1(T-\widehat\pi)^2=P_1(u-e_\pi)^2\).  Conditional on
\(\mathcal F_2\), the function \(e_\pi\) is fixed and the variables
\[
 W_i=(u_i-e_\pi(X_i))^2=(T_i-\widehat\pi(X_i))^2
\]
are conditionally iid.  Their conditional center is
\[
 m_\pi:=\E(W_i\mid\mathcal F_2)
 =v+\norm{e_\pi}_2^2,
 \tag{B.26}
\]
where sample-split independence and \(\E(u\mid X)=0\) were used,
with \(v=\E u^2\).  Put \(B_D:=B_0+B_{\mathcal G}=B_\pi\).  Since
\(0\le W_i\le B_D^2\), also
\(\abs{W_i-m_\pi}\le B_D^2\), and
\[
 \operatorname{Var}(W_i\mid\mathcal F_2)
 \le\E(W_i^2\mid\mathcal F_2)\le B_D^2m_\pi.
\]
Lemma~\ref{lem:bernstein}, conditionally on \(\mathcal F_2\), therefore
uses \(B=B_D^2\), \(\sigma^2=B_D^2m_\pi\), and defines the event
\[
 \mathcal E_{R,z}^{\rm den}:=
 \left\{\abs{D-m_\pi}\le
 B_D\sqrt{\frac{2m_\pi z}{n}}+\frac{2B_D^2z}{3n}\right\},
 \qquad
 \Pp\{(\mathcal E_{R,z}^{\rm den})^c\mid\mathcal F_2\}\le2e^{-z}.
\]
The same failure bound holds unconditionally.  With
\(c_{\rm den}:=v_-/(32B_D^2)\), the square-root term is at most
\(m_\pi/4\), and the linear term is at most
\(v_-/48\le m_\pi/48\), whenever
\(1\le z\le c_{\rm den}n\).  Consequently, on this event,
\begin{equation}
 v_-/2\le m_\pi/2\le D\le K_{\rm den}:=B_D^2,
 \qquad 1\le z\le c_{\rm den}n.
 \tag{B.27}
\end{equation}
The upper bound is also immediate from \(D=P_1W\le B_D^2\).  Thus the
floor \((v_-/2)\vee D\) is inactive on
\(\mathcal E_{R,z}^{\rm den}\) and defines the estimator on its complement.

The raw R estimator is
\[
 \widehat\beta_R^{\rm raw}
 =\frac{P_1\{(Y-\widehat m)(T-\widehat\pi)\}}
 {(v_-/2)\vee P_1(T-\widehat\pi)^2}.
\]
Define its clipped version by
\[
 \widehat\beta_R=\Pi_{[-3,3]}(\widehat\beta_R^{\rm raw}).
\]
On (B.27), direct substitution gives the exact signed identity
\begin{align}
 D(\widehat\beta_R^{\rm raw}-\beta_0)
 ={}&P_1\{\varepsilon(u-e_\pi)\}-P_1\{r_m(u-e_\pi)\}\nonumber\\
 &+\beta_0\{P_1(ue_\pi)-P_1e_\pi^2\}.
\tag{B.28}
\end{align}
This R identity is derived independently, not by symmetry: it follows from
direct substitution into the R ratio after the denominator event, and it
contains no O-specific object.
With \(r=\bar t\), \(d=d_{\rm all}\), \(w_z=\sqrt{z/n}\), and
\(\ell=\sqrt{\tau_n/\lambda_R}\), (B.9a), (B.24), and
\(\delta_z=K(r+w_z)\) give
\[
 \norm{e_\pi}_{1,n}\le K(y+r+w_z),\qquad
 \norm{r_m}_{1,n}\le
 K\{d+r+w_z+(d/r)w_z+\ell\}.
\]
Choose a deterministic represented \(g_\pi\), available by
Lemma~\ref{lem:sequential-density-transfer}, such that
\(\norm{g_\pi-\pi_0}_2\le b+\rho_n\).  Its fixed-comparator event is
\[
 \mathcal E_{R,z}^{\rm cmp}:=
 \{\norm{g_\pi-\pi_0}_{1,n}\le K_{R,\rm cmp}(b+\rho_n+w_z)\}.
\]
Bounded scalar Bernstein applied to \((g_\pi-\pi_0)^2\) gives
\(\Pp\{(\mathcal E_{R,z}^{\rm cmp})^c\}\le2e^{-z}\).  Define also
\[
\begin{split}
 \mathcal E_{R,z}^{\rm wt}:=\bigg\{&
 \abs{P_1\{\varepsilon(u-e_\pi)\}}
 \le K_{R,\rm wt}\norm{u-e_\pi}_{1,n} w_z,\\
 &\abs{P_1(r_m u)}\le K_{R,\rm wt}\norm{r_m}_{1,n} w_z,\\
 &\abs{P_1(ue_\pi)}\le K_{R,\rm wt}\norm{e_\pi}_{1,n} w_z\bigg\}.
\end{split}
\]
The first inequality follows from (B.1a) conditional on
\(\mathcal F_T\); the other two follow from (B.1a) conditional on
\(\mathcal F_X\).  Taking expectations of each conditional failure over
its own conditioning field and then applying an unconditional union bound
gives failure at most \(6e^{-z}\).  The remaining fitted propensity calibration is
supplied by (B.25).  Decomposing
\(e_\pi=(\widehat\pi-g_\pi)+(g_\pi-\pi_0)\), (B.1a), (B.25), and empirical
Cauchy--Schwarz yield
\begin{align*}
 \abs{P_1\{\varepsilon(u-e_\pi)\}}
  &\le K(1+y+r+w_z)w_z,\\
 \abs{P_1(r_mu)}
  &\le K\{d+r+w_z+(d/r)w_z+\ell\}w_z,\\
 \abs{P_1(r_me_\pi)}
  &\le K\{(y+r+w_z)^2+(r+w_z)^2
       +\lambda_R(d+r+w_z)^2+\tau_n\\
  &\hspace{22mm} +(b+w_z)[d+r+w_z+(d/r)w_z+\ell]\},\\
 \abs{P_1(ue_\pi)}+P_1e_\pi^2
  &\le K\{(y+r+w_z)w_z+(y+r+w_z)^2\}.
\end{align*}
On the finite event
\[
 \mathcal E_{R,z}:=
 \mathcal E_z^{\rm proc}\cap\mathcal E_{\pi,z}^{\rm tr,1}\cap
 \mathcal E_{\mu,z}^{\rm tr,1}\cap\mathcal E_{R,z}^{\rm den}\cap
 \mathcal E_{R,z}^{\rm cmp}\cap\mathcal E_{R,z}^{\rm wt},
 \qquad \Pp(\mathcal E_{R,z}^c)\le K_{R,0}e^{-z},
\]
the uniform events precede every fitted substitution, so no random D1
index enters a pointwise event.  Adding the four displays and dividing by
the floor \(D\ge v_-/2\) gives the complete tail for the raw output.
Projection cannot increase distance to \(\beta_0\in[-3,3]\), so the same
tail holds for the clipped \(\widehat\beta_R\):
\begin{align}
 \abs{\widehat\beta_R-\beta_0}
 \le K_{12}\bigg[&bd+y^2+r^2+\lambda_Rd^2+\tau_n+\ell(b+r)\nonumber\\
 &+(1+b+d+y+r)w_z+\frac{bd}{r}w_z
 +\left(1+\frac d r\right)w_z^2+w_z^3+\ell w_z\bigg],
 \tag{B.29}
\end{align}
Choose \(c_3\le c_{\rm den}\wedge1\) and
\(K_{13}\ge K_{R,0}\).  Then (B.29) has failure at most
\(K_{13}e^{-z}\), for \(1\le z\le c_3n\).  The possible amplification is
explicit.
\begin{lemma}[R all-\(z\) absorption certificate]
\label{lem:r-all-z-absorption}
On the R-safe cell, the complete stochastic accounting for (B.29) is
\[
\begin{array}{c|c|c}
\text{tail monomial}&\text{layer-cake scale}&\text{deterministic control}\\ \hline
(1+b+d+y+r)w_z&q&Kq\\
\frac{bd}{r}w_z&\frac{bd}{r}q&K(q+bx+y^2)\\
w_z^2&q^2&Kq\\
\frac d r w_z^2&\frac d{nr}&Kq\\
w_z^3&q^3&Kq\\
\ell w_z&\ell q&K(q+bx+y^2).
\end{array}
\]
The deterministic terms \(bd,y^2,r^2,\lambda_Rd^2,\tau_n\), and
\(\ell(b+r)\) obey the same bound, including zero budgets, both cap
branches, and cap equality.  The events are pointwise in \(z\), not
simultaneous in \(z\); unsafe cells use the public fallback.
\end{lemma}
\begin{proof}
From
\[
 r\ge C_{\rm rad}h_n,\quad d\le B_{\rm all},\quad
 d\le K_d(x+y+r+\rho_n),\quad \lambda_Rd^2=K_\lambda r^2,
\]
we get
\[
 bd\le K(bx+y^2+q),\qquad
 \frac{bd}{r}q\le K(bx+y^2+q),\qquad
 \frac d{nr}\le\frac{B_{\rm all}}
 {C_{\rm rad}\sqrt{n\log(en)}}\le Kq.
\]
Indeed, \(h_n\le x\) sends \(bh_n\) to \(bx\), whereas
\(x<h_n\) and \(b\le y<x\) send it to \(h_n^2\le2q\).
Also \(r^2\le K(y^2+q)\) and \(\ell=\rho_nd/r\), so the same split
controls \(\ell(b+r)+\ell q\).  The other rows follow from fixed envelopes
and \(q^2,q^3\le q\).  The floor \(r=C_{\rm rad}(t\vee h_n)>0\) covers
\(t=0\), and the two bounds on \(d\) cover both cap branches and equality.
The layer-cake moments are the same gamma integrals as in
Lemma~\ref{lem:o-all-z-absorption}.  Below \(z=1\) use the \(z=1\)
threshold; above \(c_3n\) use clipping; if \(c_3n<1\), then
\(q>\sqrt{c_3}\) absorbs the bounded loss.  No simultaneous-in-\(z\) event
is asserted.
\end{proof}
Lemma~\ref{lem:r-all-z-absorption} and (B.29) prove
\begin{equation}
 \sup_{P\in\cP(\Gamma)}
 \E_P\abs{\widehat\beta_R-\beta(P)}
 \le C_R\{q+bx+y^2\}.
 \tag{B.30}
\end{equation}

\subsection{Proof-event ledger}
\label{app:upper-event-ledger}

For auditability, the following is the complete event ledger used above.
Each item records the definition, conditioning field, random variables,
failure bound, admissible \(z\)-range, and downstream use.  The ledger is a
cross-reference, not an additional intersection or a claim of simultaneity
over \(z\).

\begin{enumerate}
\item \(\mathsf P_{\pi,2}^{u}(z)\) is the simultaneous (B.3)--(B.4)
event from Lemma~\ref{lem:localized-process} for the represented treatment
difference core on \(\cD_2\), with random multiplier \(u\).  It is
unconditional, concerns the iid \(\cD_2\) observations, has failure
\(K_{\pi,2}^{u}e^{-z}\) for \(1\le z\le n\), and is used in
(B.6)--(B.7) and the treatment empirical-radius conversion.

\item \(\mathsf P_{\mu,2}^{\rm sq}(z)\) is the (B.4) event for the
represented outcome difference core on \(\cD_2\).  It is unconditional,
concerns the \(\cD_2\) covariates, has failure
\(K_{\mu,2}^{\rm sq}e^{-z}\) for \(1\le z\le n\), and is used in the
outcome empirical-radius conversion following (B.9).

\item \(\mathsf P_{J,2}^{\varepsilon}(z)\) is the simultaneous
(B.3)--(B.4) event for \(\mathcal J_\mu\) on \(\cD_2\), with multiplier
\(\varepsilon\).  It is unconditional, concerns the \((X,T,Y)\) variables
in \(\cD_2\), has failure \(K_{J,2}^{\varepsilon}e^{-z}\) for
\(1\le z\le n\), and is used in (B.8)--(B.9).

\item \(\mathsf P_{\mu,1}^{u}(z)\) is the simultaneous (B.3)--(B.4)
event for the represented outcome difference core on \(\cD_1\), with
multiplier \(u\).  It is unconditional, concerns the iid \(\cD_1\)
observations, has failure \(K_{\mu,1}^{u}e^{-z}\) for
\(1\le z\le n\), and its square part is used in (B.9b), while its
multiplier and square parts are used in (B.12)--(B.16).

\item \(\mathsf P_{\pi,2}^{\xi}(z)\) is the (B.3) event for the
represented treatment difference core on \(\cD_2\), with
\(\xi=Y-m_0(X)\).  It is unconditional, concerns \(\cD_2\), has failure
\(K_{\pi,2}^{\xi}e^{-z}\) for \(1\le z\le n\), and is used in
(B.21)--(B.22).

\item \(\mathsf P_{\pi,2}^{m_0}(z)\) is the centered (B.3) event for
the same represented core on \(\cD_2\), with fixed multiplier \(m_0(X)\).
It is unconditional, concerns \(\cD_2\), has failure
\(K_{\pi,2}^{m_0}e^{-z}\) for \(1\le z\le n\), and is used in
(B.21)--(B.22).

\item \(\mathsf P_{\pi,1}^{m_0}(z)\) is the simultaneous fixed-multiplier
(B.3) and square (B.4) event for the represented treatment difference core
on \(\cD_1\).  It is unconditional, concerns \(\cD_1\), has failure
\(K_{\pi,1}^{m_0}e^{-z}\) for \(1\le z\le n\), and is used in
 (B.9a) and (B.21)--(B.25).

\item \(\mathcal E_z^{\rm proc}\) is exactly the finite intersection of
the preceding seven primitive events.  It is unconditional, involves both
sample halves as just itemized, satisfies (B.5) for \(1\le z\le n\), and
is used by the initial fits, O calibration, and R calibration.  It contains
none of the scalar, comparator, denominator, or fitted weighted events.

\item For each fixed positive proof slack,
\(\mathcal E_{\pi,z}^{\rm cmp}(\varepsilon_{\rm cmp,\pi})\) is the
displayed one-sided bound on \(P_2r_\pi^2\).  It is unconditional, involves
only the fixed comparator residual on \(\cD_2\), has failure at most
\(2e^{-z}\) for \(1\le z\le n\), and is used in (B.6)--(B.7).
The exact treatment-fit event \(\mathcal E_{\pi,z}^{\rm fit,2}\) is then
defined by (B.7) and its \(\cD_2\) empirical analogue after scalar
tail-probability continuity; it is unconditional, concerns
\(\widehat\pi\) and \(\cD_2\), and has failure
\(K_{\pi,\rm fit,2}e^{-z}\) for \(1\le z\le n\).

\item For each fixed positive proof slack,
\(\mathcal E_{\mu,z}^{\rm cmp}(\varepsilon_{\rm cmp,\mu})\) is the
displayed one-sided bound on \(P_2(e^\circ)^2\).  It is unconditional,
involves the fixed joint comparator on \(\cD_2\), has failure at most
\(2e^{-z}\) for \(1\le z\le n\), and is used in (B.8)--(B.9).
The exact joint-fit event \(\mathcal E_{\mu,z}^{\rm fit,2}\) is defined by
(B.9) and its \(\cD_2\) empirical analogue after scalar tail-probability
continuity; it is unconditional, concerns
\((\widehat\beta_0,\widehat\mu)\) and \(\cD_2\), and has failure
\(K_{\mu,\rm fit,2}e^{-z}\) for \(1\le z\le n\).

\item For each fixed positive slack, the D1 fixed-residual events
\(\mathcal E_{\pi,z}^{\rm cmp,1}(\varepsilon_{\rm cmp,\pi})\) and
\(\mathcal E_{\mu,z}^{\rm cmp,1}(\varepsilon_{\rm cmp,\mu})\) are scalar
Bernstein events for \(g_\pi-\pi_0\) and \(g^\circ-\mu_0\) on
\(\cD_1\), each with failure at most \(2e^{-z}\).  Conditional on
\(\mathcal F_2\), the fitted represented indices are fixed, while the D1
square-process events are uniform before insertion.  Combining these facts
with (B.7)--(B.9), then using scalar tail-probability continuity rather than
an infinite event intersection, defines
\(\mathcal E_{\pi,z}^{\rm tr,1}\) and
\(\mathcal E_{\mu,z}^{\rm tr,1}\).  They are unconditional, have failures
\(K_{\pi,\rm tr,1}e^{-z}\) and \(K_{\mu,\rm tr,1}e^{-z}\) for
\(1\le z\le n\), and supply (B.9a)--(B.9b) at every D1 use site:
(B.13), (B.18)--(B.20), (B.23)--(B.25), and (B.29)--(B.30).

\item \(\mathcal E_{D_0,z}\) is the displayed scalar Bernstein event for
\(D_0=P_1(Tu)\).  It is unconditional, involves \((T,u)\) on \(\cD_1\),
has failure at most \(2e^{-z}\), and for \(1\le z\le c_0n\) supplies the
normalization and denominator bounds used in (B.12)--(B.18).

\item \(\mathcal E_{O,z}^{\rm cmp}\) is the displayed empirical-norm
event for the fixed \(r_\mu=\mu_0-g_\mu\).  It is unconditional, involves
the \(\cD_1\) covariates, has failure at most \(2e^{-z}\) for
\(1\le z\le n\), and is used in the third and fourth terms of (B.17).
The event
\(\mathcal E_{O,z}^{\rm wt}\) is the displayed intersection of three
weighted inequalities: the first two are conditional on \(\mathcal F_T\)
and the third on \(\mathcal F_X\); they involve respectively
\(\varepsilon w\), \(\varepsilon(\widehat a-w)\), and \(r_\mu w\).
After integrating each conditional bound, their unconditional union failure
is at most \(6e^{-z}\) for \(1\le z\le n\); they are used in
(B.17)--(B.18).

\item \(\mathcal E_{O,z}\) is exactly the displayed six-event finite
intersection of \(\mathcal E_z^{\rm proc}\), the two D1 transfer events,
\(\mathcal E_{D_0,z}\), and the O comparator and weighted events.  It is
unconditional after integrating the conditional bounds, has failure
\(K_{O,0}e^{-z}\) for \(1\le z\le c_2n\), and supports (B.18), whose
layer-cake integration gives (B.20).

\item \(\mathcal E_{R,z}^{\rm den}\) is the displayed conditional
Bernstein event for \(W_i=(T_i-\widehat\pi(X_i))^2\), conditional on
\(\mathcal F_2\).  Because \(e_\pi\) is \(\mathcal F_2\)-measurable,
conditioning fixes it and leaves the \(\cD_1\) variables \(W_i\)
conditionally iid, with center
\(m_\pi=v+\norm{e_\pi}_2^2\), envelope \(0\le W_i\le B_D^2\), and
variance proxy
\(\operatorname{Var}(W_i\mid\mathcal F_2)\le B_D^2m_\pi\).  The event
has conditional and unconditional failure at most \(2e^{-z}\) for
\(1\le z\le c_{\rm den}n\); the choices
\(c_{\rm den}=v_-/(32B_D^2)\) and \(K_{\rm den}=B_D^2\) imply (B.27) and
supply the denominator in (B.28)--(B.29).

\item \(\mathcal E_{R,z}^{\rm cmp}\) is the displayed empirical-norm
event for the fixed \(g_\pi-\pi_0\).  It is unconditional, involves the
\(\cD_1\) covariates, has failure at most \(2e^{-z}\) for
\(1\le z\le n\), and is used after the decomposition following (B.28).
The event
\(\mathcal E_{R,z}^{\rm wt}\) is the displayed intersection of three
weighted inequalities: the first is conditional on \(\mathcal F_T\) and
the other two on \(\mathcal F_X\); they involve
\(\varepsilon(u-e_\pi)\), \(r_m u\), and \(ue_\pi\).  After integrating
each conditional bound, their unconditional union failure is at most
\(6e^{-z}\) for \(1\le z\le n\); they are used in (B.29).

\item \(\mathcal E_{R,z}\) is exactly the displayed six-event finite
intersection of \(\mathcal E_z^{\rm proc}\), the two D1 transfer events,
and the R denominator, comparator, and weighted events.  It is unconditional
after integrating the conditional bounds, has failure
\(K_{R,0}e^{-z}\) for \(1\le z\le c_3n\), and supports (B.29), whose
layer-cake integration gives (B.30).
\end{enumerate}

\subsection{Public four-cell selector and Q1 lift}
\label{app:upper-selector}

Define \(T_\mu=ay+x^2\), \(T_\pi=bx+y^2\), and
\[
 I_O=\ind\{x\le y,\ \bar s\le\vartheta_\rho\},\qquad
 I_R=\ind\{y<x,\ \bar t\le\vartheta_\rho\}.
\]
Define
\[
 \widehat\beta_\Gamma=I_O\widehat\beta_O+I_R\widehat\beta_R.
\]
When both indicators vanish the value is zero.  Ties go to O and safe
equalities remain safe.  Since the indicators are disjoint and zero belongs
to \([-3,3]\), branchwise clipping equals clipping after selection.

Let \(W=T_\mu\wedge T_\pi\).  If \(x\le y\), then
\(T_\mu\le2xy\le2T_\pi\); the mirror holds when \(y<x\).  Put
\(V=\min\{a+s^2,xy,T_\mu,T_\pi\}\).  Directly,
\[
 V\le W\le4V,\qquad X/5\le V\le6X.
 \tag{B.31}
\]
For the last inequality, use the following six exhaustive weak-order cells,
with the stated boundary convention:
\[
\begin{array}{c|c}
 (a\ge t,b\le s)&T_\pi\le6X\\
 (a\le t,b\ge s)&T_\mu\le6X\\
 (a\ge t,b\ge s,a\ge b)&T_\pi\le6X\\
 (a\ge t,b\ge s,a<b)&T_\mu\le6X\\
 (a\le t,b\le s,s\le t)&T_\mu\le6X\\
 (a\le t,b\le s,t<s)&T_\pi\le6X.
\end{array}
\]
Weak equalities are assigned by the displayed first matching row, so the
cells cover all zeros and ties without changing the bound.  This proves
\begin{equation}
 T_{\rm selected}\le2W\le48X.
 \tag{B.32}
\end{equation}

For the zero fallback use the already frozen \(\delta_\rho\).  If
\(h_n>\delta_\rho\), then \(\delta_\rho^2<h_n^2\le2q\).  If
\(h_n\le\delta_\rho\) and the O-oriented cell is unsafe, then
\(s>\delta_\rho\).  Either \(t\ge s/2\) and \(M\ge s^2/4\), or
\(t<s/2\) and \(x\le y\) imply \(P>s^2/2\).  The R-oriented argument is
the exact mirror.  Since clipped zero loses at most three, on every unsafe
cell
\begin{equation}
 \abs{0-\beta(P)}\le12\delta_\rho^{-2}(q+X)
 \tag{B.33}
\end{equation}
on every unsafe cell.

Equations (B.20), (B.30), (B.32), and (B.33) prove, uniformly for every
fixed legal \(\Gamma\),
\[
 \sup_{P\in\cP(\Gamma)}
 \E_P\abs{\widehat\beta_\Gamma-\beta(P)}\le C_\rho(q+X).
\]
Every fit, edit, indicator, fallback, and clipping operation is Borel by the
countable and finite-dimensional constructions above.  Therefore this rule
is a legal competitor inside the fixed-\(\Gamma\) infimum.  Taking the outer
supremum only after the uniform fixed-pair inequality proves the upper half
of Theorem~\ref{thm:main}; no \(\sup/\inf\) exchange and no cross-pair
estimator is used.


\subsection{Proof of the allocation corollary}
\label{app:allocation-proof}
\begin{proof}
Theorem~\ref{thm:main} and \(W\le T_{\rm selected}\le2W\), \(X/5\le W\le24X\), give \(\risk\asymp_\rho q+T_{\rm selected}\). Division by \(xy>0\) gives the first assertion, including necessity because all three terms are nonnegative.
For the allocation bounds, \(X\le4xy\): each of \(AB,M,O,P\) is at most \(xy\). Also \(X\ge x^2/4\): one of \(a,s\) is at least \(x/2\), and one of \(b,t\) is at least \(y/2\); the four combinations force respectively \(ab,O,P,M\ge x^2/4\). Finally the allocations \((0,x,y,0)\) and \((x,0,y,0)\) have \(X=x^2(1+y)\) and \(X=xy\), respectively.
\end{proof}
